\label{der:7-3}
设$z_1,\ldots,z_k$是$U\cap W$的基，分别扩充成$U$的基$(z_1,\ldots,z_k,u_1,\ldots,u_a)$及$W$的基$(z_1,\ldots,z_k,w_1,\ldots,w_b)$。拼接后生成$U+W$。若有零关系，则$\sum_j c_jw_j$等于属于$U$的向量，同时属于$W$，所以在交空间中；用交基表示后，与$W$的基无关性比较，得全$c_j=0$；剩余关系在$U$基中，所有其他系数也为零。因此拼接为基，维数为$k+a+b=\dim U+\dim W-k$。

对商空间，先取$U$的基$u_1,\ldots,u_k$并补$v_1,\ldots,v_{n-k}$成$V$基。任意向量的陪集只保留后面$v_j$的组合，所以$v_j+U$生成商空间。若$\sum_j c_j(v_j+U)=U$，则$\sum_jc_jv_j\in U$；与完整基比较得$c_j=0$。故这些陪集为基，商维数$n-k$。陪集运算与代表元无关，因为换代表元只增加$U$内的向量，其和与数乘仍在$U$中。

对偶方面，对$v=\sum_jx_jb_j$定义$f_i(v)=x_i$。基坐标唯一，使函数良定义且线性，满足$f_i(b_j)=\delta_{ij}$。任意线性函数$f$有$f(v)=\sum_jx_j f(b_j)=\sum_j f(b_j)f_j(v)$，所以$f=\sum_jf(b_j)f_j$；对零关系作用于$b_i$，得系数全零，因此$f_i$构成$n$个元素的对偶基，$\dim V^*=n$。
